{\centering{\bfseries{ВВЕДЕНИЕ}}\par}
\addcontentsline{toc}{section}{ВВЕДЕНИЕ}

В эпоху цифровой трансформации данные стали ключевым ресурсом для анализа, машинного обучения и принятия решений.
Однако использование реальных данных, особенно в чувствительных областях (медицина, финансы, персональная аналитика), сталкивается с проблемами конфиденциальности и законодательных ограничений, таких как GDPR и ФЗ-152 «О персональных данных».

Цель выпускной квалификационной работы --- разработка метода обеспечения дифференциальной приватности при генерации синтетических данных. 

Для достижения поставленной цели необходимо решить следующие задачи:
\begin{itemize}[label=---]
	\item проанализировать и классифицировать существующие методы генерации синтетических данных и дифференциальной приватности;
	\item сформулировать концепцию комбинированного метода на уровне идеи, формализовать задачу в виде функциональной модели;
	\item выделить ключевые этапы разрабатываемого метода и описать их в виде детализированной диаграммы IDEF0 и схем алгоритмов;
	\item разработать программное обеспечение, реализующее разработанный метод;
	\item исследовать характеристики разработанного метода при различных входных
	данных;
	\item провести сравнение результатов работы реализованного метода с альтернативными методами.
	
\end{itemize}

\clearpage